Los datos obtenidos del horno eléctrico se muestran en \cref{fig:DataHorno}, donde se tomaron un total de 109 mediciones de temperatura en intervalos de 20 segundos, hasta que se alcanzó una estabilidad en la temperatura del horno. 

\begin{figure}[htbp!]
    \centering
    \includegraphics[width=0.9\linewidth]{Image/data-plots_B.pdf}
    \caption{Datos obtenidos de un horno eléctrico}
    \label{fig:DataHorno}
\end{figure}

En \cref{fig:FirstOrden} se muestran los resultados de los modelos de primer orden. Se desarrolló un modelo utilizando el software Matlab ($M_{f}$), además de otros dos modelos, uno basado en métodos gráficos ($H_{f_1}$) y otro utilizando el método de Alfaro ($H_{f_2}$)

\begin{figure}[htbp!]
    \centering
    \includegraphics[width=0.9\linewidth]{Image/fit-plot-firstOrder.pdf}
    \caption{Modelos de primer orden con retardo}
    \label{fig:FirstOrden}
\end{figure}

Para el modelo de segundo orden, también se realizó un modelo utilizando Matlab ($M_{s}$) y otros dos modelos empleando los métodos de $x$ ($H_{s_1}$) y $y$ ( $H_{s_2}$), como se muestra en \cref{fig:SecondOrden}. Finalmente, en la \cref{fig:IEAP} se puede observar la integral del error absoluto de predicción (IEAP) para cada uno de los metodos empleados.

\begin{figure}[htbp!]
    \centering
    \includegraphics[width=0.9\linewidth]{Image/fit-plots-SecondOrder.pdf}
    \caption{Modelos de segundo orden con retardo}
    \label{fig:SecondOrden}
\end{figure}

\begin{figure}[htbp!]
    \centering
    \includegraphics[width=0.9\linewidth]{Image/fit-plots-ieap.pdf}
    \caption{Integral del error absoluto de predicción}
    \label{fig:IEAP}
\end{figure}

En \cref{tab:Metodos} se presentan los coeficientes obtenidos para las funciones de transferencia correspondientes a los diferentes métodos empleados para los modelos de primer y segundo orden, así como las integrales del error de predicción absoluto (IEAP) y cuadrático (IECP) y el porcentaje de error de predicción (PEP). Se concluye que el método que mejor se ajusta al sistema es el de \(x\), con un IECP de \(x\), mientras que el modelo que presenta el mayor desajuste corresponde al de \(x\).

Por otro lado, se observa un alto valor de IEAP en los modelos obtenidos mediante métodos manuales, debido a que el tiempo de recolección de datos no fue suficiente para determinar experimentalmente la ganancia de la función de transferencia. En su lugar, se utilizó el valor proporcionado por Matlab.



\begin{table}[htbp!]
    \centering
    \rowcolors{2}{white}{gray!25}  % Alternancia de colores
    \begin{tabular}{c|S[table-format=1.3]|S[table-format=3.3]|S[table-format=3.3]|S[table-format=2.3]|S[table-format=4.3]|S[table-format=4.3]|S[table-format=3.3]}
    \toprule
        {Método} & {$k_{p}$} & {$Tp_{1}$} & {$Tp_{2}$} & {$Td$} & {IEAP} & {IECP} & {PEP} \\
    \midrule
    \bottomrule
    \end{tabular}
    \caption{Resultados obtenidos para los métodos empleados}
    \label{tab:Metodos}
\end{table}
